\section{The Contemporary Doctrinal Synthesis}

\subsection{The Catechism's ordered plurality}

The \emph{Catechism of the Catholic Church} 1950--1986 offers the clearest contemporary map. All law finds its first and final truth in eternal law. The expressions of moral law are interrelated: eternal law; natural law; revealed Old and New Law; and civil and ecclesiastical laws (1950--1953; \sourceurl{https://www.vatican.va/content/catechism/en/part_three/section_one/chapter_three/article_1/the_moral_law.html}{``The Moral Law''}). The list is coordinated rather than competitive.

Natural law expresses the original moral sense by which reason discerns good and evil; it is universal, stable in first principles, and the basis of fundamental rights and duties. Its application varies, and its conclusions can be obscured by sin and cultural conditions (1954--1960; \sourceurl{https://www.vatican.va/content/catechism/en/part_three/section_one/chapter_three/article_1/i_the_natural_moral_law.html}{``The Natural Moral Law''}). Revelation does not make nature unintelligible; it gives fallen humanity needed light and grace.

The Old Law prepares for Christ and includes moral, ceremonial, and juridical precepts; the New Law fulfills and surpasses it and is principally the grace of the Holy Spirit, expressed in Christ's teaching and the apostolic catechesis (1961--1974). Human civil and ecclesiastical laws belong within this economy, not beside it as sovereign rivals.

\subsection{\emph{Veritatis splendor}: participation, universality, and embodiment}

John Paul II's \emph{Veritatis splendor} 35--53 responds to theories that detach freedom from truth or oppose nature and reason. Human reason participates in eternal law; natural law is not a biological code imposed from outside but rational participation by the embodied person (\sourceurl{https://www.vatican.va/content/john-paul-ii/en/encyclicals/documents/hf_jp-ii_enc_06081993_veritatis-splendor.html}{text}). Universality does not erase singularity. Practical judgment applies stable moral truth to concrete acts; the sustained treatment at nn.~79--83 explains why intention or circumstances cannot transform an intrinsically evil moral object into a good choice.

This teaching protects both poles. Objectivity does not mean that every disputed policy is directly deduced from a first principle. Contextual judgment does not mean that the moral object is created by preference. Positive law needs prudence, and prudence is not permission to redefine the good.

\subsection{Natural law as a public language}

The International Theological Commission's 2009 \emph{In Search of a Universal Ethic} surveys wisdom traditions, Greco-Roman sources, Scripture, Christian developments, moral experience, first principles, historicity, political society, and Christological fulfillment (\sourceurl{https://www.vatican.va/roman_curia/congregations/cfaith/cti_documents/rc_con_cfaith_doc_20090520_legge-naturale_en.html}{text}). It is an expert theological commission document, not an act of the Magisterium with the authority of an encyclical. Its value is methodological: natural-law reasoning seeks common moral intelligibility rather than using ``nature'' as an unexplained confessional password.

Public reason need not exclude theological motivation from citizens. \emph{Gaudium et spes} 43 and 76 preserve both Christian responsibility and legitimate plurality in temporal judgments. The same natural principle can support more than one reasonable positive specification. Catholic social teaching supplies ends and principles---dignity, solidarity, subsidiarity, common good, preferential concern for the vulnerable---without becoming a complete municipal code.

\subsection{Church teaching and Church legislation}

The teaching office and the power of governance are related but not identical. The Magisterium authoritatively teaches faith and morals; a legislator enacts juridical norms; an executive authority applies law; a tribunal decides controversies. One papal document can exercise more than one function, but its genre, words of enactment, approval, promulgation, subjects, and effective date determine the legal effect.

This yields four recurring patterns:

\begin{itemize}
\item A doctrinal document may teach a divine norm without creating a new canonical penalty.
\item A canon may declare a divine norm and add human procedural protections.
\item A disciplinary law may rest on theological reasons while remaining mutable ecclesiastical positive law.
\item A curial explanation may be authoritative guidance without being an authentic interpretation having the force of law.
\end{itemize}

The 1983 CIC itself distinguishes these functions. Canon 747 \S2 claims the Church's competence to announce moral principles concerning the social order and judge human affairs insofar as fundamental rights or salvation require it; canons 749--754 calibrate forms of teaching and assent (\sourceurl{https://www.vatican.va/archive/cod-iuris-canonici/eng/documents/cic_lib3-cann747-755_en.html}{CIC 747--755}). Those canons should not be read as if every prudential ecclesiastical statement were a universal legislative command.

\subsection{Natural right and positive recognition}

A positive constitution can recognize a right grounded in human dignity. The state does not create the person's dignity, yet its legal recognition is not redundant: it supplies jurisdiction, standing, remedies, procedures, institutional duties, and limits on officials. The right is natural in ground and positive in juridical form.

\emph{Dignitatis humanae} illustrates the pattern. Religious immunity is grounded in the nature and dignity of the person and is to be recognized in constitutional law as a civil right. Positive recognition makes a naturally grounded claim legally actionable. Limits must themselves be ordered by objective moral norms and just public order, not majority discomfort.

\subsection{The limits of conscience language}

Conscience does not confer immunity from every external consequence, and law cannot simply compel interior faith. \emph{Dignitatis humanae} distinguishes the moral duty to seek truth from civil immunity against coercion. \emph{Evangelium vitae} distinguishes resistance to gravely unjust commands from the prudential analysis of cooperation and harm reduction. A legal system may establish procedures for conscientious objection; their scope is a positive-law question informed by, but not identical to, the moral claim.

For a concrete case, one must identify the act required, moral object, degree of participation, causal and institutional setting, available alternatives, rights of others, competent exemption process, and consequences in each forum. ``My conscience says'' and ``the law says'' are both incomplete without that analysis.

\subsection{What remains disputed}

Catholic thinkers disagree about the best philosophical formulation of natural law, the relation of basic goods to metaphysics, the status of \emph{ius gentium}, the scope of legal validity language, and many prudential applications. Thomistic, Augustinian, personalist, new-natural-law, virtue-jurisprudential, and classical juridical approaches overlap without becoming identical. The Church's teaching does not settle every intramural philosophical dispute.

The study's taxonomy is therefore a governing grammar, not a machine. It identifies burden of proof and competent questions. It does not remove the need for exegesis, history, moral theology, empirical knowledge, legislation, or adjudication.
